% !TeX program = xelatex
\documentclass[a4paper,10pt]{article}

% ============================================================================
% សន្លឹកកិច្ចការសិស្ស ស-ខ្មៅ (Student Worksheet - Grade 8 2.4.4 POW10)
% បរិបទ៖ គ្មានស្លាយ LCD និងគ្មានសម្ភារៈកាត់បិទ — សិស្សបំពេញផ្ទាល់លើសន្លឹកកិច្ចការ ២ ទំព័រ A4
% ============================================================================
\usepackage[a4paper, left=0.85cm, right=0.85cm, top=0.70cm, bottom=0.95cm, footskip=0.48cm]{geometry}
\usepackage{amsmath, amssymb}
\usepackage{unicode-math}
\usepackage{fontspec}
\usepackage{array, colortbl}
\usepackage[table]{xcolor}
\usepackage{enumitem}
\usepackage{pifont}
\usepackage{fancyhdr}
\usepackage{lastpage}
\usepackage{tikz}
\usepackage[most]{tcolorbox}
\usetikzlibrary{calc, arrows.meta, decorations.pathreplacing, positioning, babel}

\makeatletter
\let\@expl@@@mark@update@singlecol@structures@@\relax
\makeatother

\providecommand{\NGSFindFonts}{%
  \IfFileExists{C:/Windows/Fonts/cambria.ttc}{%
    \def\NGSSysPath{C:/Windows/Fonts/}%
    \def\NGSUserPath{C:/Users/KIT/AppData/Local/Microsoft/Windows/Fonts/}%
  }{%
    \IfFileExists{./fonts/KhmerOSsiemreap.ttf}{\def\NGSSysPath{./fonts/}\def\NGSUserPath{./fonts/}}{%
    \IfFileExists{../fonts/KhmerOSsiemreap.ttf}{\def\NGSSysPath{../fonts/}\def\NGSUserPath{../fonts/}}{%
    \IfFileExists{../../fonts/KhmerOSsiemreap.ttf}{\def\NGSSysPath{../../fonts/}\def\NGSUserPath{../../fonts/}}{%
    \IfFileExists{../../../fonts/KhmerOSsiemreap.ttf}{\def\NGSSysPath{../../../fonts/}\def\NGSUserPath{../../../fonts/}}{%
    \IfFileExists{../../../../fonts/KhmerOSsiemreap.ttf}{\def\NGSSysPath{../../../../fonts/}\def\NGSUserPath{../../../../fonts/}}{%
    \IfFileExists{../../../../../fonts/KhmerOSsiemreap.ttf}{\def\NGSSysPath{../../../../../fonts/}\def\NGSUserPath{../../../../../fonts/}}{%
    \IfFileExists{../../../../../../fonts/KhmerOSsiemreap.ttf}{\def\NGSSysPath{../../../../../../fonts/}\def\NGSUserPath{../../../../../../fonts/}}{%
      \def\NGSSysPath{../../../../../../../fonts/}\def\NGSUserPath{../../../../../../../fonts/}%
    }}}}}}}%
  }%
}
\NGSFindFonts

\setmathfont[Path={\NGSSysPath}, FontIndex=1, Scale=1.02]{cambria.ttc}
\defaultfontfeatures{Ligatures=TeX}

\setmainfont[
  Path={\NGSUserPath}, Extension=.ttf, Script=Khmer,
  BoldFont={KhmerOSsiemreap}, BoldFeatures={FakeBold=1.8},
  ItalicFont={KhmerOSsiemreap}, ItalicFeatures={FakeSlant=0.22}
]{KhmerOSsiemreap}

\newfontfamily\muol[Path={\NGSUserPath}, Extension=.ttf, Script=Khmer]{KhmerOSmuollight}
\newfontfamily\TimesEn[Path={\NGSSysPath}, Extension=.ttf, BoldFont={timesbd}, ItalicFont={timesi}]{times}

\XeTeXlinebreaklocale "km"
\XeTeXlinebreakskip = 0pt plus 1pt minus 0.1pt

\definecolor{panelgray}{gray}{0.93}
\definecolor{headgray}{gray}{0.85}

\pagestyle{fancy}
\fancyhf{}
\renewcommand{\headrulewidth}{0pt}
\fancyfoot[L]{\fontsize{7.5pt}{9.5pt}\selectfont អនុវិទ្យាល័យព្រះក្រសាំង \en{|} គណិតវិទ្យាថ្នាក់ទី៨ \en{|} មេរៀនទី២៖ ៤.៤. ស្វ័យគុណនៃ \en{10} (គ្រូ៖ ហុក ប៊ុនថែន)}
\fancyfoot[R]{\fontsize{7.8pt}{9.5pt}\selectfont ទំព័រ \en{\thepage/\pageref{LastPage}}}

\setlength{\parindent}{0pt}
\setlength{\parskip}{0pt}
\setlength{\tabcolsep}{3.8pt}

\newcommand{\en}[1]{{\TimesEn #1}}
\newcommand{\enbf}[1]{{\TimesEn\bfseries #1}}
\newcommand{\unit}[1]{\,\text{\en{#1}}}
\newcommand{\cnum}[1]{\mbox{\fontsize{10pt}{12pt}\selectfont\ding{\numexpr181+#1\relax}}}
\newcommand{\chkbox}{\tikz[baseline=-0.12ex]\draw[black, line width=0.8pt] (0,0) rectangle (0.24cm,0.24cm);}
\newcommand{\blank}[1][1.5cm]{\makebox[#1]{\dotfill}}
\newcommand{\ansbox}[1][0.85cm]{\,\tikz[baseline=-0.5ex]\draw[black, line width=0.65pt, fill=white] (0,0) rectangle (#1, 0.42cm);\,}
\newcommand{\secbar}[2]{%
  \begin{tcolorbox}[enhanced, sharp corners, colback=headgray, colframe=black, boxrule=0.9pt, boxsep=1pt, top=2pt, bottom=2pt, left=5pt, right=5pt, before skip=4pt, after skip=3pt]
    {\muol\fontsize{9.2pt}{13pt}\selectfont #1}\hfill{\fontsize{8pt}{11pt}\selectfont\textbf{#2}}
  \end{tcolorbox}%
}

\newcolumntype{L}[1]{>{\raggedright\arraybackslash}p{#1}}
\newcolumntype{C}[1]{>{\centering\arraybackslash}p{#1}}

\begin{document}
\fontsize{9.6pt}{14.4pt}\selectfont

% ============================================================================
% ទំព័រទី ១៖ E1 + E2 (POE 1 & POE 2) + E3 សំយោគ និងធនាគារហេតុផល
% ============================================================================
\begin{tcolorbox}[enhanced, sharp corners, colback=white, colframe=black, boxrule=1.0pt, boxsep=2pt, top=3pt, bottom=3pt, left=6pt, right=6pt]
  \begin{center}
    {\muol\fontsize{10.5pt}{15pt}\selectfont អនុវិទ្យាល័យព្រះក្រសាំង \enbf{--} សន្លឹកកិច្ចការសិស្ស \enbf{(}មេរៀនទី២៖ ៤\enbf{.}៤\enbf{.} ស្វ័យគុណនៃ \enbf{10)}}\\[3pt]
    {\fontsize{8.8pt}{12pt}\selectfont
    ឈ្មោះសិស្ស / ដៃគូ៖ \blank[5.2cm] \hspace{0.25cm}
    ថ្នាក់ទី៨៖ \blank[1.1cm] \hspace{0.25cm}
    ក្រុមទី៖ \blank[1.0cm] \hspace{0.25cm}
    កាលបរិច្ឆេទ៖ ១៧/១១/២០២៦}
  \end{center}
\end{tcolorbox}

% --- E1 + E2 ---
\secbar{ផ្នែកទី១៖ បេសកកម្មទស្សន៍ទាយ \enbf{(E1)} និងការរុករកជាដៃគូ \enbf{(E2 -- POE 1 \& POE 2)}}{រយៈពេល ១៤ នាទី}

\begin{tcolorbox}[enhanced, sharp corners, colback=panelgray, colframe=black, boxrule=0.75pt, boxsep=1.5pt, top=2.5pt, bottom=2.5pt, left=5pt, right=5pt]
  \textbf{បេសកកម្ម \en{E1} (គូសធីក \chkbox\ មុនពិសោធន៍)៖}
  (១) តម្លៃដឹកស្រូវ $3.75\times 10^4\text{ ៛}$ ស្មើ៖ \chkbox\ $3.750000$ \quad \chkbox\ $37\,500$ \quad
  \en{|} \quad (២) ចំណុះថ្នាំ $10^{-3}\unit{l}$ ស្មើ៖ \chkbox\ $-1000\unit{l}$ \quad \chkbox\ $0.001\unit{l}$
\end{tcolorbox}

\vspace{2pt}
\textbf{ក. សកម្មភាពទី ១ (\en{POE 1})៖ ជណ្ដើរឆ្លុះស្វ័យគុណនៃ $10$ (សៀវភៅពុម្ព ទំ.១៩–២០ — ចែកនឹង $10$ ម្ដងមួយកាំ)}

{\renewcommand{\arraystretch}{1.34}
\begin{tabular}{|C{1.90cm}|C{3.35cm}|C{2.45cm}||C{2.00cm}|C{3.25cm}|C{4.65cm}|}
  \hline
  \rowcolor{headgray}
  \textbf{ស្វ័យគុណ $+$} & \textbf{តម្លៃលេខ} & \textbf{ចំនួនសូន្យ} & \textbf{ស្វ័យគុណ $-$} & \textbf{ប្រភាគ $\dfrac{1}{10^n}$} & \textbf{ទសភាគ និងចំនួនខ្ទង់ក្រោយក្បៀស} \\
  \hline
  $10^0$ & $1$ & $0$ សូន្យ & $10^0$ & $\dfrac{1}{1}$ & $1$ (ចំណុចកណ្ដាលជណ្ដើរឆ្លុះ!) \\
  \hline
  $10^1$ & $10$ & $1$ សូន្យ & $10^{-1}$ & $\dfrac{1}{10}$ & $0.1$ ($1$ ខ្ទង់ក្រោយក្បៀស) \\
  \hline
  $10^2$ & $100$ & $2$ សូន្យ & $10^{-2}$ & $\dfrac{1}{100}$ & $0.01$ ($2$ ខ្ទង់ក្រោយក្បៀស) \\
  \hline
  $10^3$ & $1\,000$ & $\blank[0.8cm]$ សូន្យ & $10^{-3}$ & $\dfrac{1}{\blank[1.2cm]}$ & $\blank[1.5cm]$ ($\blank[0.6cm]$ ខ្ទង់ក្រោយក្បៀស) \\
  \hline
  $10^4$ & $\blank[2.0cm]$ & $\blank[0.8cm]$ សូន្យ & $10^{-4}$ & $\dfrac{1}{10\,000}$ & $\blank[1.5cm]$ ($\blank[0.6cm]$ ខ្ទង់ក្រោយក្បៀស) \\
  \hline
  $10^5$ & $100\,000$ & $\blank[0.8cm]$ សូន្យ & $10^{-5}$ & $\dfrac{1}{100\,000}$ & $\blank[1.6cm]$ ($\blank[0.6cm]$ ខ្ទង់) \\
  \hline
  $10^6$ & $\blank[2.2cm]$ & $\blank[0.8cm]$ សូន្យ & $10^{-6}$ & $\dfrac{1}{10^6}$ & $\blank[1.8cm]$ ($\blank[0.6cm]$ ខ្ទង់) \\
  \hline
  $10^7$ & $10\,000\,000$ & $\blank[0.8cm]$ សូន្យ & $10^{-7}$ & $\dfrac{1}{10^7}$ & $0.\blank[2.0cm]$ ($\blank[0.6cm]$ ខ្ទង់) \\
  \hline
\end{tabular}}

\vspace{3pt}
$\bullet$ \textbf{ពន្យល់ (\en{E})៖} ពេលចែកនឹង $10$ តម្លៃកាន់តែតូចជិតសូន្យ ដូច្នេះ $10^{-3} = 0.001$ ជាចំនួន\blank[2.0cm] (មិនមែនអវិជ្ជមានទេ!) ។

\vspace{4pt}
\textbf{ខ. សកម្មភាពទី ២ (\en{POE 2})៖ ផ្លូវធ្នូលោតក្បៀសពេលគុណ ឬចែកនឹង $10^n$ (ឧទាហរណ៍ ទំ.២០)}

{\renewcommand{\arraystretch}{1.42}
\begin{tabular}{|L{3.55cm}|L{6.25cm}|L{4.75cm}|C{3.25cm}|}
  \hline
  \rowcolor{headgray}
  \textbf{ប្រមាណវិធី} & \textbf{គណនាតាមការបំបែកតម្លៃ} & \textbf{ទិសដៅ និងចំនួនខ្ទង់លោតក្បៀស} & \textbf{លទ្ធផលត្រឹមត្រូវ} \\
  \hline
  \textbf{ក.} $3.75 \times 10^4$ & $= 3.75 \times 10\,000 = 3\,\underset{\smile}{7}\,\underset{\smile}{5}\,\underset{\smile}{0}\,\underset{\smile}{0}.$ & លោតទៅ\textbf{ស្ដាំ} $\blank[0.7cm]$ ខ្ទង់ (ថែមសូន្យ $\blank[0.6cm]$) & $= \blank[1.8cm]$ \\
  \hline
  \textbf{ខ.} $128.5 \div 10^2$ & $= \dfrac{128.5}{100} = 128.5 \times 10^{-2} = 1.\underset{\smile}{2}\,\underset{\smile}{8}\,5$ & លោតទៅ\textbf{ឆ្វេង} $\blank[0.7cm]$ ខ្ទង់ & $= \blank[1.8cm]$ \\
  \hline
  \textbf{គ.} $4.2 \times 10^{-3}$ & $= 4.2 \div 10^3 = 0.\underset{\smile}{0}\,\underset{\smile}{0}\,\underset{\smile}{4}\,2$ & លោតទៅ\blank[1.2cm] $\blank[0.7cm]$ ខ្ទង់ & $= \blank[1.8cm]$ \\
  \hline
  \textbf{ឃ.} $0.09 \times 10^4$ & $= 0.09 \times 10\,000$ & លោតទៅ\blank[1.2cm] $\blank[0.7cm]$ ខ្ទង់ & $= \blank[1.8cm]$ \\
  \hline
\end{tabular}}

% --- E3 ---
\secbar{ផ្នែកទី២៖ សំយោគក្បួនទូទៅ \enbf{(E3)} និងធនាគារហេតុផល \enbf{B\&W} \cnum{1}\enbf{--}\cnum{5}}{សៀវភៅពុម្ព ទំ.១៩–២០}

\begin{tabular}{|L{9.20cm}|L{9.20cm}|}
  \hline
  \textbf{១. ក្បួនជាទូទៅ (ទំ.១៩–២០)៖}\par\vspace{3pt}
  $\bullet$ ចំពោះ $n$ ជាចំនួនគត់វិជ្ជមាន ($n\in\mathbb{Z}^+$)៖\par\vspace{2pt}
  $10^n = 1\underbrace{00\dots 0}_{\text{សូន្យ }\blank[0.6cm]}$ \quad និង \quad $10^{-n} = \dfrac{1}{10^n} = 0.\underbrace{00\dots 01}_{\blank[0.6cm]\text{ ខ្ទង់}}$\par\vspace{3pt}
  $\bullet$ ដើម្បី\textbf{គុណ}ចំនួនទសភាគនឹង $10^n$ ត្រូវរំកិលក្បៀសទៅខាង\blank[1.3cm] ចំនួន $\blank[0.6cm]$ ខ្ទង់ ។\par\vspace{2pt}
  $\bullet$ ដើម្បី\textbf{ចែក}ចំនួនទសភាគនឹង $10^n$ (ឬគុណនឹង $10^{-n}$) ត្រូវរំកិលក្បៀសទៅខាង\blank[1.3cm] ចំនួន $\blank[0.6cm]$ ខ្ទង់ ។
  &
  \textbf{២. ធនាគារហេតុផល \en{B\&W} (\cnum{1}--\cnum{5})៖}\par\vspace{2pt}
  \cnum{1} $10^n = 1\underbrace{00\dots 0}_{n\text{ សូន្យ}}$ ($n\in\mathbb{Z}^+$) និង $10^0 = 1$ ។\par\vspace{2pt}
  \cnum{2} $10^{-n} = \frac{1}{10^n} = 0.\underbrace{00\dots 01}_{n\text{ ខ្ទង់}} > 0$ (មិនមែនចំនួនអវិជ្ជមានទេ!) ។\par\vspace{2pt}
  \cnum{3} $\times 10^n$ ($n>0$) $\to$ រំកិលក្បៀសទៅ\textbf{ស្ដាំ} $n$ ខ្ទង់ (ថែម $0$ បើខ្វះខ្ទង់) ។\par\vspace{2pt}
  \cnum{4} $\div 10^n$ ឬ $\times 10^{-n} \to$ រំកិលក្បៀសទៅ\textbf{ឆ្វេង} $n$ ខ្ទង់ ។\par\vspace{2pt}
  \cnum{5} ប្ដូរខ្នាត៖ $1\unit{cm}=10^{-2}\unit{m}$ ($\div 100$), \; $1\unit{g}=10^{-3}\unit{kg}$ ($\div 1000$), \; $1\unit{ml}=10^{-3}\unit{l}$ ($\div 1000$) ។ \\
  \hline
\end{tabular}

\vspace{4pt}
{\renewcommand{\arraystretch}{1.38}
\begin{tabular}{|C{4.40cm}|C{4.40cm}|C{4.40cm}|C{4.40cm}|}
  \hline
  \rowcolor{panelgray}
  $10^5 = \blank[2.0cm]$ & $10^{-5} = \blank[2.0cm]$ & $0.0001 = 10^{\blank[0.7cm]}$ & $1\,000\,000 = 10^{\blank[0.7cm]}$ \\
  \hline
  $0.01 = 10^{\blank[0.7cm]}$ & $10^0 = \blank[1.5cm]$ & $5.4\times 10^2 = \blank[1.5cm]$ & $9\div 10^3 = \blank[1.6cm]$ \\
  \hline
\end{tabular}}

\newpage

% ============================================================================
% ទំព័រទី ២៖ E4 លំហាត់គំរូ ១ & ២ + E5 ប្រតិបត្តិ ១ & ២ + វិនិច្ឆ័យ A-D + Exit Check + កិច្ចការផ្ទះ
% ============================================================================
\secbar{ផ្នែកទី៣៖ ពង្រីកចំណេះដឹងតាមជំហានស៊ើបអង្កេត \enbf{(E4 --} លំហាត់គំរូ ទំ\enbf{.}២០\enbf{)}}{រយៈពេល ១០ នាទី}

\textbf{១. លំហាត់គំរូទី ១ (ទំ.២០)៖} គូសធ្នូលោតក្បៀស និងគណនាផលគុណ-ផលចែកខាងក្រោម៖

{\renewcommand{\arraystretch}{1.40}
\begin{tabular}{|L{3.55cm}|L{8.95cm}|C{3.35cm}|C{1.95cm}|}
  \hline
  \rowcolor{headgray}
  \textbf{លំហាត់} & \textbf{ការវិភាគទិសដៅ និងចំនួនខ្ទង់លោតក្បៀស} & \textbf{ចម្លើយ} & \textbf{ហេតុផល} \\
  \hline
  \textbf{ក.} $62.013 \times 10^5$ & លោតទៅស្ដាំ $5$ ខ្ទង់៖ ឆ្លង $013$ (បីខ្ទង់) រួចថែមសូន្យ $\blank[0.8cm]$ ខ្ទង់ទៀត & $= \blank[2.3cm]$ & \cnum{3} \\
  \hline
  \textbf{ខ.} $0.57 \div 10^5$ & $= 0.57\times 10^{-5} \to$ លោតទៅឆ្វេង $5$ ខ្ទង់ (ថែមសូន្យ $\blank[0.8cm]$ ពីមុខលេខ $57$) & $= \blank[2.3cm]$ & \cnum{4} \\
  \hline
\end{tabular}}

\vspace{4pt}
\textbf{២. លំហាត់គំរូទី ២ (ទំ.២០)៖} ចម្លង និងបំពេញតារាងប្ដូរខ្នាតប្រវែងពីសង់ទីម៉ែត្រ ($\unit{cm}$) ទៅម៉ែត្រ ($\unit{m}$) ($\div 10^2 = \div 100$)៖

{\renewcommand{\arraystretch}{1.42}
\begin{tabular}{|L{4.80cm}|C{4.20cm}|C{4.20cm}|C{4.35cm}|}
  \hline
  \rowcolor{headgray}
  \textbf{ប្រវែងគិតជាសង់ទីម៉ែត្រ ($\unit{cm}$)} & $10\unit{cm}$ & $25\unit{cm}$ & $136\unit{cm}$ \\
  \hline
  ប្រវែងគិតជាម៉ែត្រ ($\unit{m}$) \textbf{ជាប្រភាគសម្រួល} & $\dfrac{10}{100}\unit{m} = \dfrac{1}{10}\unit{m}$ & $\dfrac{25}{100}\unit{m} = \blank[1.6cm]\unit{m}$ & $\dfrac{136}{100}\unit{m} = \blank[1.6cm]\unit{m}$ \\
  \hline
  ប្រវែងគិតជាម៉ែត្រ ($\unit{m}$) \textbf{ជាទសភាគ} & $0.1\unit{m}$ & $\blank[1.8cm]\unit{m}$ & $\blank[1.8cm]\unit{m}$ \\
  \hline
\end{tabular}}

\vspace{4pt}
$\bullet$ \textbf{សំណួរឆ្លុះបញ្ចាំង \en{(EIA Looking-Back)}៖} ហេតុអ្វី $0.57\div 10^5 = 0.0000057$ មានសូន្យក្រោយក្បៀស $5$ ខ្ទង់ តែ $34.57\div 10^5$ មានសូន្យក្រោយក្បៀសតែ $3$ ខ្ទង់? ចម្លើយ៖ \blank[5.8cm]

% --- E5 ---
\secbar{ផ្នែកទី៤៖ វាយតម្លៃលទ្ធផលសិក្សា \enbf{(E5 --} ប្រតិបត្តិ ទំ\enbf{.}២១ \enbf{+} វិនិច្ឆ័យ \enbf{A--D +} សំបុត្រចេញ\enbf{)}}{រយៈពេល ១០ នាទី}

\textbf{ក. លំហាត់ប្រតិបត្តិ ១ និង ២ (សៀវភៅពុម្ព ទំ.២១ — ធ្វើជាដៃគូ រួចប្ដូរសន្លឹកផ្ទៀង \en{Pause-Review})៖}\\[2pt]
\textbf{១. គណនាផលគុណ និងផលចែកខាងក្រោម៖}\hspace{0.4cm}
\textbf{ក.} $762.0213 \times 10^5 = \blank[2.8cm]$ \hspace{0.5cm} \en{|} \hspace{0.5cm}
\textbf{ខ.} $34.57 \div 10^5 = \blank[2.8cm]$

\vspace{3pt}
\textbf{២. បំពេញតារាងប្ដូរខ្នាតទម្ងន់ពីក្រាម ($\unit{g}$) ទៅគីឡូក្រាម ($\unit{kg}$) ជាប្រភាគសម្រួល និងទសភាគ ($1\unit{kg}=10^3\unit{g}=1000\unit{g}$)៖}

{\renewcommand{\arraystretch}{1.42}
\begin{tabular}{|L{5.10cm}|C{4.10cm}|C{4.10cm}|C{4.25cm}|}
  \hline
  \rowcolor{headgray}
  \textbf{ទម្ងន់គិតជាក្រាម ($\unit{g}$)} & $10\unit{g}$ & $325\unit{g}$ & $1365\unit{g}$ \\
  \hline
  ទម្ងន់គិតជាគីឡូក្រាម ($\unit{kg}$) \textbf{ជាប្រភាគសម្រួល} & $\dfrac{10}{1000} = \blank[1.5cm]\unit{kg}$ & $\dfrac{325}{1000} = \blank[1.5cm]\unit{kg}$ & $\dfrac{1365}{1000} = \blank[1.5cm]\unit{kg}$ \\
  \hline
  ទម្ងន់គិតជាគីឡូក្រាម ($\unit{kg}$) \textbf{ជាទសភាគ} & $\blank[1.8cm]\unit{kg}$ & $\blank[1.8cm]\unit{kg}$ & $\blank[1.8cm]\unit{kg}$ \\
  \hline
\end{tabular}}

\vspace{4pt}
\textbf{ខ. គូសវង់ជុំវិញជម្រើសត្រឹមត្រូវតែ ១ គត់ ក្នុងចំណោម \en{A, B, C, D} រួចសរសេរលេខហេតុផល \cnum{1}--\cnum{5} នៃជម្រើសខុស៖}

{\renewcommand{\arraystretch}{1.30}
\begin{tabular}{|C{4.40cm}|C{4.40cm}|C{4.40cm}|C{4.40cm}|}
  \hline
  \cellcolor{panelgray}\textbf{ជម្រើស \en{A}} & \cellcolor{panelgray}\textbf{ជម្រើស \en{B}} & \cellcolor{panelgray}\textbf{ជម្រើស \en{C}} & \cellcolor{panelgray}\textbf{ជម្រើស \en{D}} \\
  \hline
  $0.57\div 10^5 = 0.0000057$ & $3.75\times 10^4 = 3.750000$ & $10^{-4} = -10\,000$ & $10\unit{g} = 0.1\unit{kg}$ \\[3pt]
  \chkbox\ ត្រូវ \quad \chkbox\ ខុស (\blank[0.7cm]) & \chkbox\ ត្រូវ \quad \chkbox\ ខុស (\blank[0.7cm]) & \chkbox\ ត្រូវ \quad \chkbox\ ខុស (\blank[0.7cm]) & \chkbox\ ត្រូវ \quad \chkbox\ ខុស (\blank[0.7cm]) \\
  \hline
\end{tabular}}

\vspace{4pt}
\begin{tcolorbox}[enhanced, sharp corners, colback=white, colframe=black, boxrule=0.85pt, boxsep=2pt, top=3pt, bottom=3pt, left=5pt, right=5pt]
  {\muol\fontsize{8.8pt}{12.5pt}\selectfont សំបុត្រចេញរហ័ស \enbf{(Exit Check --} ៦០ វិនាទី\enbf{)}៖}\\[2pt]
  (១) $10^{-5} = \blank[1.8cm]$ \quad \en{|} \quad
  (២) $4.8 \times 10^3 = \blank[1.8cm]$ \quad \en{|} \quad
  (៣) $250\unit{ml} = \blank[1.3cm]\unit{l}$ (ប្រភាគ) $= \blank[1.3cm]\unit{l}$ (ទសភាគ)
\end{tcolorbox}

\vspace{4pt}
\begin{tcolorbox}[enhanced, sharp corners, colback=panelgray, colframe=black, boxrule=0.85pt, boxsep=2pt, top=3pt, bottom=3pt, left=5pt, right=5pt]
  {\muol\fontsize{8.8pt}{12.5pt}\selectfont កិច្ចការផ្ទះ \enbf{(}សៀវភៅពុម្ព ទំព័រ ២៦ លំហាត់ \enbf{\#10} \enbf{--} ប្ដូរ $\unit{ml}\to\unit{l}$ ដោយចែកនឹង $1000$\enbf{)}៖}\\[3pt]
  {\renewcommand{\arraystretch}{1.32}\footnotesize
  \begin{tabular}{|L{3.50cm}|C{2.15cm}|C{2.15cm}|C{2.25cm}|C{2.15cm}|C{2.25cm}|C{2.35cm}|}
    \hline
    \textbf{ចំណុះ ($\unit{ml}$)} & $10\unit{ml}$ & $330\unit{ml}$ & $1365\unit{ml}$ & $100\unit{ml}$ & $3800\unit{ml}$ & $17500\unit{ml}$ \\
    \hline
    \textbf{លីត្រ ($\unit{l}$) ប្រភាគ} & $\blank[1.4cm]$ & $\blank[1.4cm]$ & $\blank[1.4cm]$ & $\blank[1.4cm]$ & $\blank[1.4cm]$ & $\blank[1.4cm]$ \\
    \hline
    \textbf{លីត្រ ($\unit{l}$) ទសភាគ} & $\blank[1.4cm]$ & $\blank[1.4cm]$ & $\blank[1.4cm]$ & $\blank[1.4cm]$ & $\blank[1.4cm]$ & $\blank[1.4cm]$ \\
    \hline
  \end{tabular}}
\end{tcolorbox}

\end{document}
